% ===========================================================================
% Chapter: Transporting lift tests into MMM parameters (dive 02)
% ===========================================================================
\chapter{Transporting lift tests into MMM parameters}
\label{ch:transport}

\begin{provenance}
\textbf{Source dives:} \dive{02} \emph{A formal transport map from lift tests to MMM parameters (M2)} (run 2026-08-31, file \texttt{02\_experiment\_transport\_map.md}, last written 2026-08-31 23:14 UTC).
\textbf{Code:} \texttt{code/02\_experiment\_transport\_map.py}.
\textbf{Frontier-study problem(s) addressed:} M2 --- write the Pearl--Bareinboim selection diagram for experiment-to-MMM calibration, derive when the experimental functional identifies the model functional, characterise the correction terms (adstock truncation, saturation-point mismatch), and handle several experiments as a joint, correlated posterior over the ROAS vector.
\textbf{Backlog items closed:} none; \textbf{opened:} \BL{3} (calibration under misspecification), \BL{4} (a specification test for transport).
\end{provenance}

\begin{keyideas}
\begin{itemize}
\item In the selection diagram for experiment-to-MMM transport the effect of spend is \emph{not} nonparametrically transportable; every calibration rests on a parametric \emph{shape-invariance} assumption, under which transport is mediated by $\theta$ and nothing licenses the scalar map iROAS $\to$ model ROI.
\item Exact identity $\iROAS(d) = m_\infty(\theta)\,\rho\,C\,Q$ with closed-form window capture $C = 1 - \alpha^{P+1}(1-\alpha^{T_e})/(T_e(1-\alpha))$ (verified to $9\times10^{-7}$): a 4-week test on an $\alpha = 0.8$ channel analysed to test end observes 41\% of the sales it caused.
\item Naive calibration is \emph{precisely wrong}: iROAS as a prior on marginal ROAS gives $-43.0\%$ median bias (theory $-43.8\%$) at the tightest spread of any mode (IQR 0.016), and flips to $+67.3\%$ when the test runs in geos at half the national spend.
\item The \emph{operator} likelihood $\hat L \sim \mathcal N(\mathcal E_d(\theta), \mathrm{se}^2)$, which simulates the experiment's design inside the model, needs no correction factors, halves the uncalibrated curve error, wrecks no fits and has $65\times$ lower median decision loss than the ROI-prior mode.
\item A scalar iROAS carries rank-1 Fisher information; the per-period trajectory of the same experiment carries rank 4 (standalone SE 0.20 on $\alpha$); joint correlated posteriors over many experiments come free from entering each as a likelihood row on the shared $\theta$.
\end{itemize}
\end{keyideas}

\section{Problem and motivation}

An incrementality experiment measures a finite-horizon intention-to-treat effect of a specific spend change, at a specific operating point on the response curve, in a specific subpopulation, over a specific window. An MMM parameter---an average or marginal return on spend---is a steady-state structural quantity for the target population. Calibration practice equates the two: the experiment's incremental ROAS (iROAS) and its standard error become the mean and standard deviation of a prior on the model's ROI. Practitioners document the symptom (``reconciliation headaches rather than trust-building calibration'' when the experiment's KPI, geography or window differ from the model's; warnings that experiment-derived priors ``transfer bias across non-comparable populations'') without the mechanism. Problem M2 asked for the mechanism: instantiate the Pearl--Bareinboim transportability calculus, derive the exact map between the experimental and model functionals, characterise the correction terms, and say how several experiments should enter one posterior.

\dive{02} answers all four parts in the single-channel adstock/Hill model of \cref{ch:identification}: a negative theorem, an exact three-factor identity with a closed form for its largest factor, a Monte Carlo demonstration that the standard shortcut is biased by tens of percent while \emph{appearing} more certain than any alternative, and a constructive fix---calibrate on the experiment's design, not on a renamed scalar---that is the marketing form of indirect inference. The dive predates the program's attack-and-refine protocol (see the reviewer's note).

\section{Background: the ideas this chapter builds on}

\subsection{Selection diagrams and transportability (Pearl and Bareinboim, 2011, 2014, 2016)}

Transportability asks when a causal relation learned in a source population $\Pi$ (where an experiment ran) can be carried to a target $\Pi^\star$ (where only observational data exist). Starting from the causal diagram $G$ the two share, a \emph{selection diagram} adds a selection node $S_i \to V_i$ for every variable whose generating mechanism may differ between populations; the absence of an $S$-node asserts invariance. Writing $P^\star(v) = P(v \mid s^\star)$, the target effect $P^\star(y \mid do(x))$ is \emph{transportable} if it is uniquely computable from the source experimental distribution $P(y \mid do(x))$ and target observational data $P^\star(v)$. The simplest sufficient condition is \emph{s-admissibility}: if pre-treatment covariates $Z$ satisfy $(Y \perp S \mid X, Z)$ in $G_{\bar X}$ (arrows into $X$ removed) then
\begin{equation}
P^\star(y \mid do(x)) = \sum_z P(y \mid do(x), z)\,P^\star(z),
\label{eq:transport:formula}
\end{equation}
the \emph{transport formula}: experimental conditional effects re-weighted by the target's covariate distribution. The complete characterisation (Pearl and Bareinboim, 2011, 2014) is that $P^\star(y \mid do(x))$ is transportable if and only if $P(y \mid do(x), s)$ can be reduced by do-calculus to an expression in which $S$ appears only in do-free terms; the 2016 PNAS paper extends this to fusing many sources. The negative side matters here: if an $S$-node points directly into the outcome $Y$ and nothing observed screens it off, no derivation exists and the effect is not transportable without parametric assumptions.

\subsection{The geometric-adstock, Hill-saturation MMM (Jin et al., 2017)}

The Bayesian MMM modern packages descend from models a KPI as baseline plus, per channel, a \emph{carryover} transform followed by a \emph{shape} transform of spend $x_t$. Carryover is geometric adstock, $a_t = x_t + \alpha a_{t-1}$ (a unit of spend contributes $\alpha^l$ after $l$ periods); the shape is the Hill function
\begin{equation}
h(a; K, S) = \frac{a^S}{a^S + K^S},
\label{eq:transport:hill}
\end{equation}
with half-saturation $K$ and slope $S$, convex below its inflection $a^\dagger = K((S-1)/(S+1))^{1/S}$ and concave above; the channel contributes $\beta h(a_t)$ with $\beta$ the sales ceiling. Jin et al.\ show that with weekly national data the likelihood is flat along ridges in $(\alpha, K, S)$ and posteriors are prior-dominated (\cref{ch:identification} gives the mathematics). At steady spend $\bar x$ the adstock settles at $\bar a = \bar x/(1-\alpha)$; the business consumes the \emph{average ROAS} $\beta h(\bar a)/\bar x$ and the \emph{steady-state marginal ROAS} $m_\infty = \beta h'(\bar a)/(1-\alpha)$, the long-run sales per extra unit of sustained spend.

\subsection{What ``calibrating to a lift test'' means in practice}

\emph{Meridian} (Google) follows Zhang, Wurm et al.\ (2024): the model is reparameterised so the prior sits not on the coefficient $\beta_m$ but on the channel's ROI (incremental KPI over the data window divided by spend), $\beta_m$ becoming a deterministic function of ROI and the other parameters; a lift experiment supplies the mean and standard deviation of a lognormal ROI prior (an mROI variant, the return to a small spend increase, is also offered). The documentation concedes ``there is no single formula to translate an experiment result into a prior''. \emph{Robyn} (Meta; Runge et al., 2024) is a ridge MMM whose adstock and saturation hyperparameters come from an evolutionary search over a Pareto front; calibration enters as a third objective, the discrepancy (MAPE.LIFT) between the model's decomposed channel contribution over the experiment's date window and the experiment's absolute lift. \emph{PyMC-Marketing} holds the strongest prior art: a lift test with baseline $x$ and change $\Delta x$ enters the posterior as a Normal likelihood on the \emph{static saturation secant} $\beta[s(x + \Delta x) - s(x)]$---exact on curvature, but with no adstock dynamics, window truncation or population selection (Orduz's simulation study works in the same spirit).

\subsection{Indirect inference (Gouri\'eroux, Monfort and Renault, 1993)}

When a structural model can be simulated but its likelihood is awkward, indirect inference fixes an \emph{auxiliary} statistic $\hat\beta$, defines the \emph{binding function} $b(\theta)$ as the value that statistic takes on data simulated from the model at $\theta$, and sets $\hat\theta = \argmin_\theta (\hat\beta_{\mathrm{obs}} - b(\theta))^\T\Omega(\hat\beta_{\mathrm{obs}} - b(\theta))$. The auxiliary quantity is never \emph{equated} with a structural parameter, even one of the same name; the model is asked what value of the statistic it would produce, and that is matched. Under misspecification the estimator converges to a pseudo-true value.

\subsection{The adstock transient as a dose-ranging sweep (\cref{ch:identification}; Heusch, 2026)}

\dive{01} found that a two-level spend design identifies all four parameters through its \emph{transitions}: each step drags $a_t$ along a known exponential path, tracing an arc of the Hill curve, and discarding transient observations collapses the smallest Fisher eigenvalue about $700\times$. Heusch (2026) recovers adstock, saturation and effectiveness structurally from the temporal variation a geo experiment induces by the same fact. \cref{transport:prop:rank} applies the observation to what an experiment \emph{reports}.

\section{Setup and assumptions}

The model is that of \cref{ch:identification}: $y_t = \beta\,h(a_t; K, S) + \varepsilon_t$, $a_t = x_t + \alpha a_{t-1}$, $\varepsilon_t \sim \mathcal N(0, \sigma^2)$, $\theta = (\beta, K, S, \alpha)$, working point $\theta_0 = (1, 1.5, 2, 0.6)$, national mean spend $\bar x = 1$ (so $\bar a = 2.5$), $\sigma = 0.05$. Subscripts $\star$ and $e$ mark the target domain $\Pi_\star$ (national, business-as-usual spend) and the experiment domain $\Pi_e$ (a subpopulation over a window).

\begin{definition}[Design and functionals]
\label{transport:def:design}
A design is $d = (\bar x_e, \{\delta_t\}_{t < T_e}, W = T_e + P)$: baseline spend $\bar x_e$ in the experiment geos, a pulse $\delta_t$ for $T_e$ periods, and a window of $W$ periods extending $P$ periods past the pulse. The \emph{experiment functional} is the model-implied total incremental KPI over the window,
\begin{equation}
\mathcal E_d(\theta) = \beta_e \sum_{t \in W}\Big[h\big(a_t(x_e + \delta)\big) - h\big(a_t(x_e)\big)\Big], \qquad
\iROAS(d) = \frac{\mathcal E_d(\theta)}{\sum_t \delta_t},
\label{eq:transport:Ed}
\end{equation}
where $x_e$ is the baseline spend path (constant at $\bar x_e$) and adstock starts at its steady state $\bar x_e/(1-\alpha)$; the experiment reports an estimate $\hat L$ of $\mathcal E_d$ with standard error se. The \emph{model functionals} are $m_\infty(\theta) = \beta_\star h'(\bar a_\star)/(1-\alpha)$ and the average ROAS $\beta_\star h(\bar a_\star)/\bar x_\star$. $\Delta a_t$ and $\Delta h_t$ denote the adstock and Hill differences the pulse induces.
\end{definition}

\begin{assumption}[Shape invariance]
\label{transport:ass:shape}
$(K, S, \alpha)$ are shared between $\Pi_e$ and $\Pi_\star$; only $\beta$ (and seasonality/demand) may differ.
\end{assumption}

\begin{assumption}[Internal validity and model family]
\label{transport:ass:internal}
$\hat L$ is unbiased for the intention-to-treat effect under $do(X)$ (no contamination or spillover); adstock is geometric and, for the closed form of $C$, the pulse is uniform.
\end{assumption}

\section{Main results}

\subsection{No nonparametric transport: all calibration is parametric}

The shared graph is $X_t \to A_t \to Y_t$ with an unobserved demand state $U \to X_t$, $U \to Y_t$ (severed by $do(X)$ in $\Pi_e$), plus three selection nodes: $S_\beta \to Y$ (effectiveness differs by geo), $S_x \to X$ (operating point differs), $S_g \to Y$ (seasonality and demand differ).

\begin{proposition}[Non-transportability]
\label{transport:prop:nontransport}
\tagEstablished\ In this selection diagram the causal effect of spend on sales is not nonparametrically transportable from $\Pi_e$ to $\Pi_\star$: no do-calculus derivation of $P^\star(y \mid do(x))$ from $P(y \mid do(x))$ and $P^\star(v)$ exists.
\end{proposition}

$S_\beta$ points directly into the outcome mechanism and nothing observed screens it off, so $P(y \mid do(x), s)$ cannot be made $s$-free. The dive calls this ``direct from established theory, but worth stating sharply'' because of its consequence: every experiment-to-MMM calibration rides on a parametric invariance assumption, usually the implicit \cref{transport:ass:shape}. Under it the experiment \emph{is} informative about $\theta$ and all transport is $\theta$-mediated (experiment $\to$ likelihood on $\theta$ $\to$ any target functional). Nothing licenses a scalar map from iROAS to a model ROI; the identity below prices that shortcut.

\subsection{The transport identity}

\begin{theorem}[Transport identity]
\label{transport:thm:identity}
\tagEstablished\ For any design $d$ and any $\theta$,
\begin{equation}
\iROAS(d) = m_\infty(\theta)\cdot\underbrace{\rho}_{\text{selection}}\cdot\underbrace{C(\alpha; T_e, P)}_{\text{window capture}}\cdot\underbrace{Q(\delta, \bar a_e; \theta)}_{\text{curvature/secant}},
\label{eq:transport:identity}
\end{equation}
with $\rho = \beta_e h'(\bar a_e)/(\beta_\star h'(\bar a_\star))$, $C = (1-\alpha)\sum_{t \in W}\Delta a_t/\sum_t \delta_t$ and $Q = \sum_t \Delta h_t/(h'(\bar a_e)\sum_t \Delta a_t)$.
\end{theorem}

The identity is exact because $Q$ absorbs the remainder by definition; its content is the closed forms and the naming of three mechanisms. The derivation is two lines: $\mathcal E_d = \beta_e h'(\bar a_e)\,Q\sum_t\Delta a_t$ and $\sum_t \Delta a_t = C\sum_t\delta_t/(1-\alpha)$ by the definitions, so $\iROAS = \beta_e h'(\bar a_e)CQ/(1-\alpha) = m_\infty\rho CQ$.

\begin{lemma}[Window capture, closed form]
\label{transport:lem:C}
\tagEstablished\ For a uniform pulse of length $T_e$ and a window of $T_e + P$ periods,
\begin{equation}
C(\alpha; T_e, P) = 1 - \frac{\alpha^{P+1}\,(1-\alpha^{T_e})}{T_e\,(1-\alpha)}.
\label{eq:transport:C}
\end{equation}
\end{lemma}

$C$ is the fraction of the pulse's infinite-horizon adstock response that falls inside the window ($T_e = 1$ gives the folklore $1 - \alpha^{P+1}$); it follows by summing each impulse's in-window contribution $(1-\alpha^{W-s})/(1-\alpha)$ over $s < T_e$. The magnitudes are the practical point (\cref{transport:tab:C}): a 4-week test on an $\alpha = 0.8$ channel analysed only through test end has $C = 0.41$---it observes 41\% of the incremental sales it caused; at $\alpha = 0.9$, $C = 0.23$. Common practice (short tests, high-carryover channels, no post-period) bakes a 2.5--4$\times$ understatement of the long-run effect into the ``ground truth'' before any modelling starts.

\begin{proposition}[Curvature/secant factor]
\label{transport:prop:Q}
\tagEstablished\ To second order in the pulse,
\begin{equation}
Q \approx 1 + \frac{h''(\bar a_e)}{2h'(\bar a_e)}\cdot\frac{\sum_t \Delta a_t^2}{\sum_t \Delta a_t},
\label{eq:transport:Q}
\end{equation}
a transient-weighted generalisation of the static secant, which is the $\alpha \to 0$ case.
\end{proposition}

The sign of $Q - 1$ is that of $h''(\bar a_e)$: above the inflection $a^\dagger$ the experiment under-states the marginal slope ($Q < 1$); below it, as for a test in low-spend geos, it over-states ($Q > 1$). The expansion is accurate to under 2\% for $\delta/\bar x \le 0.25$ but fails for big tests (at $\bar a/K = 1.67$, $\delta = \bar x$: exact $Q = 0.51$ versus 0.17), so large pulses need exact evaluation. The \emph{selection factor} $\rho$ is an effectiveness ratio times an operating-point slope ratio: even with $\beta_e = \beta_\star$, a test in geos with $\bar a_e < \bar a_\star$ has $\rho > 1$, sitting on a steeper part of the curve.

\begin{example}[Three numbers a practitioner might equate]
\label{transport:ex:three}
At $\theta_0$ with a $+50\%$, 8-week test analysed to test end ($P = 0$, $\rho = 1$): $\iROAS = 0.219$, $m_\infty = 0.389$, average ROAS $= 0.735$---a spread of $3.4\times$ in a fully standard setting. By arithmetic from these figures $CQ = 0.563$, and $C = 0.816$ from \cref{eq:transport:C}, so $Q = 0.69$: truncation and secant contribute comparably.
\end{example}

\subsection{The operator fix: calibrate on the design, not on a renamed scalar}

\begin{construction}[Operator likelihood]
\label{transport:con:operator}
\tagEstablished\ (Novel for MMM; indirect inference in spirit.) Add to the MMM posterior the likelihood term
\begin{equation}
\hat L \sim \mathcal N\big(\mathcal E_d(\theta),\ \mathrm{se}^2\big),
\label{eq:transport:operator}
\end{equation}
where $\mathcal E_d(\theta)$ of \cref{eq:transport:Ed} simulates the experiment's actual design inside the model at every posterior draw: baseline path, delta path, adstock transient, window truncation, experiment-geo operating point.
\end{construction}

No correction factors are needed because nothing is transported scalar-to-scalar; $C$, $Q$ and $\rho$ are computed implicitly and exactly at each $\theta$. In indirect-inference terms the experiment's estimator is the auxiliary statistic and $\mathcal E_d$ its binding function; PyMC-Marketing's secant likelihood is the static ($\alpha = 0$, full-window, no-selection) special case. The cost is that the experiment report must carry its design, not just $(\iROAS, \mathrm{se})$: the \emph{calibration header} of \cref{transport:tab:header}, every field of which the experimenter already has. Reporting the total lift with the design attached makes the result transport-safe; reporting the trajectory upgrades a rank-1 constraint to rank 4 (\cref{transport:prop:rank}).

\begin{table}[htbp]
\centering\small
\caption{The calibration header: design metadata an experiment must ship for $\mathcal E_d(\theta)$ to be computable (\dive{02} \S4).}
\label{transport:tab:header}
\begin{tabularx}{\textwidth}{lX}
\toprule
Field & Content and purpose \\
\midrule
channel, kpi & model channel id; KPI must match the model's $y$ \\
baseline\_spend\_path & per-period spend in the experiment geos (not just the mean): fixes $\bar a_e$, $\rho$ \\
delta\_spend\_path & per-period pulse $\delta_t$: fixes $Q$ and, with the window, $C$ \\
window, post\_period\_len & start, end, $P$: fixes the truncation \\
population & geo ids or share-of-country weights, baseline spend share: fixes selection \\
estimate & $(\hat L, \mathrm{se})$ as total incremental KPI, \emph{not} iROAS \\
trajectory (optional) & per-period treated-minus-control with $\mathrm{se}_t$: rank 1 becomes rank 4 \\
\bottomrule
\end{tabularx}
\end{table}

\subsection{The scalar summary discards identifying power; joint posteriors are free}

\begin{proposition}[Rank of the experimental information]
\label{transport:prop:rank}
\tagNumerical\ At $\theta_0$, for an 8-week pulse ($\delta = 0.5$, $\bar x_e = 1$) with 6 post-periods and per-period noise 0.05 on the differenced treated-minus-control series, the Fisher information from the full trajectory has rank 4, eigenvalues $(2.7\times10^{-5},\ 0.059,\ 7.8,\ 577)$, and a standalone Cram\'er--Rao SE of 0.20 for $\alpha$; the Fisher information from the scalar iROAS of the same experiment has rank 1.
\end{proposition}

A single number constrains one direction in $\theta$-space and carries no separate information about $\alpha$; the trajectory shows ramp-up and decay (hence $\alpha$) and its transient traces an arc of the Hill curve. The dive's corollary: $k$ scalar summaries constrain $k$ directions, so scalar calibration under shape invariance needs at least three well-separated designs plus the observational level to point-identify $(\beta, K, S, \alpha)$; trajectory calibration needs one experiment.

\begin{corollary}[Joint, correlated posterior]
\label{transport:cor:joint}
\tagEstablished\ Because each experiment enters as a likelihood row \cref{eq:transport:operator} on the shared $\theta$, the posterior over any vector of derived ROAS quantities is automatically jointly correlated: the joint posterior M2 asks for is a consequence of refusing to convert experiments into independent per-channel scalar priors.
\end{corollary}

Two experiments on one channel at different deltas jointly constrain $(K, S)$ in a way two independent ROI priors cannot express---they would simply conflict, and \cref{transport:sec:numerics} shows what conflict does.

\section{Numerical evidence}
\label{transport:sec:numerics}

\emph{Experiment A} compares \cref{eq:transport:C} with brute-force simulation of an infinitesimal pulse ($\delta = 10^{-6}$) on the grid of \cref{transport:tab:C}---$\alpha \in \{0.3, 0.6, 0.8, 0.9\}$ and six $(T_e, P)$ pairs spanning $T_e = 1$ to $13$ and $P = 0$ to $26$---with maximum discrepancy $9\times10^{-7}$. \emph{Experiment B} compares exact $Q$ with \cref{eq:transport:Q} at $\bar a/K \in \{0.5, 1, 1.67, 2.5\}$ and $\delta/\bar x \in \{0.1, 0.25, 0.5, 1\}$, giving the accuracy statements after \cref{transport:prop:Q}. \emph{Experiment D} computes the Fisher matrices of \cref{transport:prop:rank} by finite differences of the model-implied 14-period trajectory and of its sum (whose variance is $W$ times the per-period variance).

\begin{table}[htbp]
\centering\small
\caption{Window-capture factor $C(\alpha; T_e, P)$ from \cref{eq:transport:C} on the Experiment A grid, where the closed form matches simulation to $9\times10^{-7}$. Columns are (pulse length $T_e$, post-period $P$) in weeks.}
\label{transport:tab:C}
\begin{tabular}{lcccccc}
\toprule
$\alpha$ & (1, 0) & (4, 0) & (8, 0) & (8, 4) & (8, 26) & (13, 8) \\
\midrule
0.3 & 0.700 & 0.894 & 0.946 & 1.000 & 1.000 & 1.000 \\
0.6 & 0.400 & 0.674 & 0.816 & 0.976 & 1.000 & 0.998 \\
0.8 & 0.200 & 0.410 & 0.584 & 0.830 & 0.999 & 0.951 \\
0.9 & 0.100 & 0.226 & 0.359 & 0.580 & 0.959 & 0.778 \\
\bottomrule
\end{tabular}
\end{table}

\emph{Experiment C} is the head-to-head of calibration modes. Each of 60 replications draws a weakly identifying national series ($T = 104$ weeks, AR(1) log-spend with about $\pm10\%$ variation, $\sigma = 0.05$) and one well-powered experiment ($\bar x_e = 1$, $\delta = 0.5$, $T_e = 8$, $P = 0$, se $= 5\%$ of the true lift). Four modes are fitted by multistart least squares, the calibration entering as one extra standardised residual: \emph{uncalibrated}; \emph{naive-avg} (Meridian-style, $\widehat{\iROAS}$ centring a prior on average ROAS $\beta h(\bar a)$); \emph{naive-marg} ($\widehat{\iROAS}$ centring a prior on $m_\infty$); \emph{operator} (\cref{eq:transport:operator}). Metrics: median relative bias and IQR of $\hat m_\infty$; median maximum error of the fitted response curve over $x \in [0.05, 2.5]$ (units of the sales ceiling); share of ``wrecked'' replications (curve error above 0.5); median profit loss from allocating at the fitted optimum (margin 3).

\begin{table}[htbp]
\centering\small
\caption{Experiment C, matched operating point: 60 Monte Carlo replications, $T = 104$, one experiment with se $= 5\%$ of lift. Medians; the dive reports no Monte Carlo standard errors.}
\label{transport:tab:modes}
\begin{tabular}{lccccc}
\toprule
Mode & $m_\infty$ median bias & $m_\infty$ IQR & med.\ max curve err. & wrecked reps & med.\ profit loss \\
\midrule
uncalibrated & $-7.4\%$ & 0.108 & 0.262 & 3\% & 0.0057 \\
naive-avg & $+4.0\%$ & 0.301 & 0.228 & 30\% & 0.1368 \\
naive-marg & $-43.0\%$ & 0.016 & 0.352 & 12\% & 0.0212 \\
operator & $-3.8\%$ & 0.093 & 0.124 & 0\% & 0.0021 \\
\bottomrule
\end{tabular}
\end{table}

Four findings follow. First, theory predicts the naive bias to within a point: with $\rho = 1$, $CQ\rho - 1 = -43.8\%$ against $-43.0\%$ observed. Second, naive calibration is \emph{precisely wrong}: naive-marg has the tightest spread of any mode (IQR 0.016, about $7\times$ tighter than the uncalibrated fit) centred 43\% off the truth---calibration bought confidence, not accuracy, the quantitative form of the adoption study's ``precision theatre'', produced here by the sophisticated method. Third, naive-avg creates likelihood \emph{conflict}, not bias: the iROAS (0.219) contradicts the observational sales level $\beta h(\bar a) = 0.735$ that 104 weeks of data pin down; the data win on the level but the fit escapes along the identification ridge, 30\% of replications end with a wrecked curve, and median profit loss is $65\times$ the operator's. This is the mechanism behind the reconciliation headaches: the two evidence sources are made to disagree by the transport error itself. Fourth, the operator dominates everywhere.

The \emph{operating-point mismatch} variant reruns the design with the test in geos at 50\% of national spend ($\bar x_e = 0.5$, $\delta = 0.25$). Naive-marg flips to $+67.3\%$ bias (theory $+72.5\%$): same pipeline, opposite sign, depending on where the test happened to run. The operator's median curve error \emph{improves} to 0.048 (from 0.124 matched): through the correct map an off-operating-point experiment is more informative, because it probes the curve where it bends; through the naive map it is more toxic.

\section{Limitations and the devil's advocate}

The dive's own counterarguments are five. \emph{Shape invariance is an assumption, not a result}: if $(K, S, \alpha)$ differ between the experiment geos and the nation the operator transports incorrectly too; unlike the naive map it makes the assumption explicit and testable (two experiments in different strata over-determine the shared shape), but hierarchical $\beta_g$ only partially absorbs $S_\beta$ and nothing absorbs an $S_K$ node without multi-stratum experiments (\BL{4}). \emph{Misspecification transfers}: if the true adstock is delayed-peak and the model geometric, the operator matches the experiment by distorting $\theta$ to pseudo-true values; naive calibration is also wrong in that case (\BL{3}). \emph{Internal validity is taken as given}: contamination, spillover and control-market disruption are upstream. \emph{The closed form for $C$ assumes geometric adstock and a uniform pulse}; non-uniform pulses need the per-impulse sum, and sales-lag dynamics beyond adstock are not modelled. \emph{Experiment C uses MAP-style point fits}, not posteriors; the ``precisely wrong'' IQR should be re-confirmed with MCMC widths, expected to be starker. No attack-and-refine round changed a result, because none was run.

\section{Discussion: impacts}

For practice the message is that the object a lift test should hand to an MMM is a \emph{design plus a total}, not a number. Every current mechanism loses something: Meridian's ROI prior transports a scalar into a parameter measuring a different quantity (\cref{transport:ex:three} shows a $3.4\times$ gap at the working point); Robyn's MAPE.LIFT compares a windowed contribution but as an objective to trade off rather than a likelihood; PyMC-Marketing's secant is the operator with $\alpha = 0$ and no window. The operator is a small change to any of them---a forward simulation of the design at each draw---and the calibration header is the interface change that enables it. The identity also lets an experimentation platform report what it has \emph{not} measured, since $C$ follows from its own window and the channel's retention.

For CMO decision-making the ``precisely wrong'' finding is the sharper lesson: naive-marg is the mode a careful analyst would prefer on the evidence in front of them---the tightest interval---and it is the most biased. Conversely, the mismatch variant turns a governance liability (tests run wherever politics allows) into an asset under the correct map.

For theory the chapter settles M2's first three parts in the parametric setting and shows the fourth to be free. Outward links: \cref{ch:voi} uses $\mathcal E_d(\theta)$ as the forward map of its value-of-information design chooser and finds the average-ROI level that Meridian-style priors chiefly move nearly worthless for fixed-budget allocation, while \dive{04} finds the rank-4 trajectory information does \emph{not} project onto the cross-channel contrasts that decision needs (decision-value index 0.29 $\to$ 0.30): the trajectory identifies the channel, not the allocation. \cref{ch:fusion} ingests experiments through the operator in a closed loop; \cref{ch:aggregation} proposes de-biasing nationally fitted $(K, S)$ before operator transport (\BL{17}); \cref{ch:panel} would transport a panel-measured memory kernel through the same likelihood (\BL{61}).

\section{Further exploration}

Of the dive's six next steps, two became backlog items: the misspecification study (\BL{3}: delayed-peak truth fitted with a geometric operator, measure the pseudo-true bias, try a sandwich or inflated-se likelihood as a patch) and the specification test (\BL{4}: two experiments in different geo strata plus a shared-shape model, with a posterior-predictive check as a formal test of shape invariance---the field has no specification test for calibration at all). The others are an \texttt{add\_experiment\_operator()} prototype against PyMC-Marketing's API; a multi-channel operator able to express the cross-channel contamination (treated-geo budget reallocations) a scalar iROAS cannot; feeding $\mathcal E_d$ to the VOI chooser (done in \cref{ch:voi}); and an analytic account of the naive-avg conflict-to-ridge-escape mechanism.

In the compiler's judgement the most promising extension is \BL{4} combined with the trajectory header: if every experiment ships its per-period trajectory, any two experiments on one channel over-determine $(K, S, \alpha)$ and their posterior-predictive discrepancy tests shape invariance at no extra experimental cost, turning the chapter's one unfalsifiable premise into a routine diagnostic. Second is the multi-channel operator, since \cref{ch:deadband} and \cref{ch:feedback} show allocation decisions live in cross-channel contrasts the single-channel identity does not reach.

\begin{reviewernote}[(compiler)]
\textbf{Protocol.} \dive{02} predates the depth standard adopted ``as of run 03'' (attack-and-refine rounds, clean-room re-derivation, uncertainty on every headline). Nothing here has been independently re-derived within the program, and \cref{transport:tab:modes} reports medians over 60 replications with no standard errors: ``30\% wrecked'' is 18 replications, ``12\%'' is 7, ``3\%'' is 2, and the operator's ``0\%'' is a bound of roughly one in sixty. ``Within a point'' ($-43.0\%$ versus $-43.8\%$) has no standard error on the median.

\textbf{The rank-4 claim is nominal.} The smallest trajectory eigenvalue, $2.7\times10^{-5}$, is $4.7\times10^{-8}$ of the largest; the code declares rank by the threshold $10^{-8}\times$ the largest eigenvalue, so the fourth direction clears it by a factor of about five, and the CRLB along that eigenvector is $\approx 1/\sqrt{2.7\times10^{-5}} \approx 190$. The honest reading is ``rank 3 strong plus one very weak direction''. The 0.20 SE for $\alpha$ is plausible (the weak direction lives in $(\beta, K, S)$, the level ridge an experiment cannot see) but comes from inverting a matrix of condition number $\sim 2\times10^{7}$ without the $s_{\min}/s_{\max}$ audit \BL{82} later made mandatory for FIM-based standard errors. The ``three well-separated designs'' corollary is a counting heuristic, not a proved statement.

\textbf{Decision relevance was later softened.} \dive{04} found trajectory readouts ``worthless'' for fixed-budget cross-channel allocation ($\kappa$ 0.29 $\to$ 0.30); \cref{transport:prop:rank} concerns identification of one channel, not decision value.

\textbf{Prior-art survey.} The current Meridian documentation (checked 2026-09-07) describes a \texttt{CalibrationBuilder} that adjusts experiment results for spend scale, duration and recency before forming the ROI prior; the dive does not mention it. It is not the design-simulating operator, but a reader should check how much of $C$ and $Q$ it implements before crediting ``never stated in the MMM calibration literature'' with full force. The indirect-inference lineage is asserted, not surveyed.

\textbf{Consistency checks passed.} The INDEX log line agrees with the writeup on every number; $65\times$ is $0.1368/0.0021 = 65.1$; ``$7\times$ tighter'' is $0.108/0.016 = 6.75$ against the uncalibrated mode ($5.8\times$ against the operator); $m_\infty = 0.389$, average ROAS 0.735 and $C = 0.816$ reproduce from $\theta_0$. The mismatch variant's ``theory $+72.5\%$'' and the operator's 0.048 appear in the text only, with no table. The log timestamp ($\sim$21:15) precedes the file's last write (23:14 UTC) by two hours without explanation.
\end{reviewernote}

\section*{Sources}
\begingroup\raggedright
\begin{enumerate}[label={[\arabic*]},leftmargin=2.5em]
\item Pearl, J. and Bareinboim, E. (2014). \emph{External Validity: From Do-Calculus to Transportability Across Populations}. Statistical Science 29(4), 579--595. \url{https://doi.org/10.1214/14-STS486}
\item Pearl, J. and Bareinboim, E. (2011). \emph{Transportability of Causal and Statistical Relations: A Formal Approach}. Proc.\ AAAI-11. \url{https://dl.acm.org/doi/10.5555/2900423.2900462}
\item Bareinboim, E. and Pearl, J. (2016). \emph{Causal Inference and the Data-Fusion Problem}. PNAS 113(27), 7345--7352. \url{https://www.pnas.org/doi/abs/10.1073/pnas.1510507113}
\item Jin, Y., Wang, Y., Sun, Y., Chan, D. and Koehler, J. (2017). \emph{Bayesian Methods for Media Mix Modeling with Carryover and Shape Effects}. Google research report. \url{https://research.google/pubs/bayesian-methods-for-media-mix-modeling-with-carryover-and-shape-effects/}
\item Zhang, Y., Wurm, M., et al. (2024). \emph{Media Mix Model Calibration With Bayesian Priors}. Google research report. \url{https://storage.googleapis.com/gweb-research2023-media/pubtools/pdf/a09f404fdc3107fafb7a52cc5af6a80e4d0fda2b.pdf}
\item Google (2024--26). \emph{Meridian documentation: calibrate treatment priors; ROI, mROI and contribution parameterizations}. \url{https://developers.google.com/meridian/docs/advanced-modeling/roi-priors-and-calibration}
\item Runge, J., Skokan, I., Zhou, G. and Pauwels, K. (2024). \emph{Packaging Up Media Mix Modeling: An Introduction to Robyn's Open-Source Approach}. arXiv:2403.14674. \url{https://arxiv.org/abs/2403.14674}
\item PyMC Labs (2024--26). \emph{PyMC-Marketing: lift-test calibration of an MMM} (documentation notebook). \url{https://www.pymc-marketing.io/en/stable/notebooks/mmm/mmm_roas.html}
\item Orduz, J. (2024). \emph{MMM and experimental calibration: a simulation study}. Blog. \url{https://juanitorduz.github.io/mmm_roas/}
\item Gouri\'eroux, C., Monfort, A. and Renault, E. (1993). \emph{Indirect Inference}. Journal of Applied Econometrics 8(S1), S85--S118.
\item Heusch, A. (2026). \emph{Structural Estimation of MMM Parameters from Geo-Experiments}. arXiv:2608.21128. \url{https://arxiv.org/abs/2608.21128}
\item Frontier study (internal): the open-questions catalogue \texttt{02\_open\_questions.md}, entry M2, and the adoption-barriers catalogue \texttt{03\_mmm\_adoption\_barriers.md}, entries O7, B5 and M5.
\item Program (internal): \dive{01}, file \texttt{01\_mmm\_identified\_set.md}, sections 3.2 and 3.5; \dive{04}, file \texttt{04\_decision\_voi\_ext.md}, finding F4.
\end{enumerate}
\endgroup
